% Options for packages loaded elsewhere
\PassOptionsToPackage{unicode}{hyperref}
\PassOptionsToPackage{hyphens}{url}
\documentclass[
  ignorenonframetext,
]{beamer}
\newif\ifbibliography
\usepackage{pgfpages}
\setbeamertemplate{caption}[numbered]
\setbeamertemplate{caption label separator}{: }
\setbeamercolor{caption name}{fg=normal text.fg}
\beamertemplatenavigationsymbolsempty
% remove section numbering
\setbeamertemplate{part page}{
  \centering
  \begin{beamercolorbox}[sep=16pt,center]{part title}
    \usebeamerfont{part title}\insertpart\par
  \end{beamercolorbox}
}
\setbeamertemplate{section page}{
  \centering
  \begin{beamercolorbox}[sep=12pt,center]{section title}
    \usebeamerfont{section title}\insertsection\par
  \end{beamercolorbox}
}
\setbeamertemplate{subsection page}{
  \centering
  \begin{beamercolorbox}[sep=8pt,center]{subsection title}
    \usebeamerfont{subsection title}\insertsubsection\par
  \end{beamercolorbox}
}
% Prevent slide breaks in the middle of a paragraph
\widowpenalties 1 10000
\raggedbottom
\AtBeginPart{
  \frame{\partpage}
}
\AtBeginSection{
  \ifbibliography
  \else
    \frame{\sectionpage}
  \fi
}
\AtBeginSubsection{
  \frame{\subsectionpage}
}
\usepackage{iftex}
\ifPDFTeX
  \usepackage[T1]{fontenc}
  \usepackage[utf8]{inputenc}
  \usepackage{textcomp} % provide euro and other symbols
\else % if luatex or xetex
  \usepackage{unicode-math} % this also loads fontspec
  \defaultfontfeatures{Scale=MatchLowercase}
  \defaultfontfeatures[\rmfamily]{Ligatures=TeX,Scale=1}
\fi
\usepackage{lmodern}
\ifPDFTeX\else
  % xetex/luatex font selection
\fi
% Use upquote if available, for straight quotes in verbatim environments
\IfFileExists{upquote.sty}{\usepackage{upquote}}{}
\IfFileExists{microtype.sty}{% use microtype if available
  \usepackage[]{microtype}
  \UseMicrotypeSet[protrusion]{basicmath} % disable protrusion for tt fonts
}{}
\makeatletter
\@ifundefined{KOMAClassName}{% if non-KOMA class
  \IfFileExists{parskip.sty}{%
    \usepackage{parskip}
  }{% else
    \setlength{\parindent}{0pt}
    \setlength{\parskip}{6pt plus 2pt minus 1pt}}
}{% if KOMA class
  \KOMAoptions{parskip=half}}
\makeatother
\setlength{\emergencystretch}{3em} % prevent overfull lines
\providecommand{\tightlist}{%
  \setlength{\itemsep}{0pt}\setlength{\parskip}{0pt}}
% Slide layout defaults for warning-clean scientific decks.
\usepackage{etoolbox}
\IfFileExists{xurl.sty}{\usepackage{xurl}}{}
\IfFileExists{seqsplit.sty}{\usepackage{seqsplit}}{\newcommand{\seqsplit}[1]{#1}}
\protected\def\breakseq#1{\seqsplit{#1}}
\protected\def\breaktt#1{\begingroup\ttfamily\seqsplit{#1}\endgroup}
\setlength{\emergencystretch}{6em}
\tolerance=5000
\hbadness=10000
\hfuzz=1pt
\setlength{\tabcolsep}{2pt}
\AtBeginEnvironment{longtable}{\tiny\renewcommand{\arraystretch}{0.86}\setlength{\tabcolsep}{1pt}}
\AtBeginEnvironment{tabular}{\tiny\renewcommand{\arraystretch}{0.86}\setlength{\tabcolsep}{1pt}}

% Natbib and cross-reference fallbacks — slides don't load natbib
% or cleveref, but combined-PDF manuscript prose may emit these
% commands. The fallback renders citations as a bracketed key list
% and cross-references as detokenized labels so slides don't fail on
% undefined control sequences or raw underscores. \providecommand is
% a no-op if packages load later.
\providecommand{\citep}[1]{[#1]}
\providecommand{\citet}[1]{#1}
\providecommand{\citealp}[1]{#1}
\providecommand{\citeauthor}[1]{#1}
\providecommand{\citeyear}[1]{#1}
\providecommand{\cref}[1]{\texttt{\detokenize{#1}}}
\providecommand{\Cref}[1]{\texttt{\detokenize{#1}}}

\newtheorem{proposition}{Proposition}
\newtheorem{hypothesis}{Hypothesis}
\usepackage{bookmark}
\IfFileExists{xurl.sty}{\usepackage{xurl}}{} % add URL line breaks if available
\urlstyle{same}
% fallback for those not using the hyperref driver hyperxmp:
\makeatletter
\@ifundefined{xmpquote}{\newcommand{\xmpquote}[1]{#1}}{}
\makeatother
\hypersetup{
  hidelinks,
  pdfcreator={LaTeX via pandoc}}

\author{\texorpdfstring{}{}}
\date{}

\begin{document}

\section{\texorpdfstring{Discussion: Implications and Community
Recommendations
\label{sec:discussion}}{Discussion: Implications and Community Recommendations }}\label{discussion-implications-and-community-recommendations}

\begin{frame}[allowframebreaks]{Relationship to Prior Development
Directions}
\protect\phantomsection\label{relationship-to-prior-development-directions}
Knight, Cordes, and Friedman \citep{knight2022fep} identified six
development directions for systematic Active Inference literature
analysis: (1) increased scope of relevant works, (2) richer annotation
schemes, (3) integration of manual and artificial contributions, (4)
transferable approaches across fields, (5) participation by diverse
contributors, and (6) updated analyses tracking the field's evolution.
This pipeline directly addresses directions 1, 2, 3, and 6: it scales
retrieval to three databases, replaces manual annotation with LLM-driven
extraction while preserving human review pathways, and produces a
pipeline designed for incremental re-execution as new literature
appears. Directions 4 and 5---cross-field transferability and community
participation---remain open and are addressed below.
\end{frame}

\begin{frame}[allowframebreaks]{Tactical and Strategic Priorities}
\protect\phantomsection\label{tactical-and-strategic-priorities}
\begin{block}{Adopt Rigorous Reporting Metadata}
\protect\phantomsection\label{adopt-rigorous-reporting-metadata}
Papers should systematically report DOIs, ORCIDs, and explicit
hypothesis commitments. Submitted preprints should forward-link to their
published versions to prevent fragmented citation subgraphs. Our
extraction pipeline prioritizes the DOI as the canonical identifier;
failing that, deduplication cascades to arXiv IDs, Semantic Scholar IDs,
and OpenAlex IDs. Broad DOI adoption would resolve the cross-source
mismatch problem, enabling higher-resolution evidence mapping.
\end{block}

\begin{block}{Explore Open Knowledge Graph Infrastructure}
\protect\phantomsection\label{explore-open-knowledge-graph-infrastructure}
We encourage the exploration of federated nanopublication server
architectures to house community-contributed assertions. This would
enable a continuously updated living literature review that incorporates
new findings as they are published. The release of nanopub-js v0.1.0
\citep{kuhn2026nanopubjs} makes browser-based creation and querying of
nanopublications practical, enabling researchers to contribute
assertions directly from web interfaces. Integrating this approach with
the Active Inference Institute's Knowledge-Engineering infrastructure
\citep{knight2022fep} could provide the standardized semantic vocabulary
necessary for rigorous cross-study comparison.
\end{block}

\begin{block}{Standardize the Ontological Lexicon}
\protect\phantomsection\label{standardize-the-ontological-lexicon}
Immediate future extraction cycles should align assertion predicates
with the formally curated Active Inference Ontology. Enforcing shared
ontological primitives across studies will accelerate the aggregation of
evidence from otherwise siloed research communities, advancing the
interoperability goal outlined by Knight et al.~\citep{knight2022fep}.
\end{block}
\end{frame}

\begin{frame}[allowframebreaks]{Empirical and Theoretical Imperatives}
\protect\phantomsection\label{empirical-and-theoretical-imperatives}
\begin{block}{Architect Unified Performance Benchmarks}
\protect\phantomsection\label{architect-unified-performance-benchmarks}
The computational tools domain (B) lacks standardized performance
benchmarks for direct comparison against deep reinforcement learning
architectures. Establishing baseline metrics analogous to standard RL
environments (e.g., OpenAI Gym) is a prerequisite for transitioning
theoretical proposals into applied systems.
\end{block}

\begin{block}{Prioritize Empirical Validation}
\protect\phantomsection\label{prioritize-empirical-validation}
Biology (C5) and Language (C3) have established theoretical frameworks
but limited empirical validation. Targeted experiments designed to test
specific FEP-derived predictions---such as demonstrating morphogenesis
as Bayesian inference or measuring active inference advantages in
language tasks---would strengthen the evidence base beyond what further
theoretical work alone can achieve.
\end{block}
\end{frame}

\begin{frame}[fragile,allowframebreaks]{Living Review Maintenance}
\protect\phantomsection\label{living-review-maintenance}
The pipeline is designed for continuous operation rather than one-time
analysis. Incremental resume capabilities (checkpoint-based assertion
extraction, merge-on-add corpus deduplication) enable periodic
re-execution as new papers are indexed. We envision a maintenance cycle
in which the pipeline is re-run quarterly, with updated hypothesis
scores and field statistics published alongside the pipeline release.
Community contributors can extend the framework by adding custom
hypothesis definitions, alternative keyword taxonomies, or
domain-specific extraction prompts---all configurable via the YAML
configuration file without modifying source code. A complementary
long-term trajectory is toward RAG-enabled access: integrating the
nanopublication knowledge graph into a Retrieval-Augmented Generation
architecture \citep{fan2024survey} would enable natural-language
querying of the evidence base, making quantitative literature synthesis
accessible to researchers without programming expertise.

\begin{block}{Agentic Workspaces and MCP Integration}
\protect\phantomsection\label{agentic-workspaces-and-mcp-integration}
Beyond traditional open-source maintenance, the repository is
architected as an intrinsically agentic workspace. Every underlying
source module (e.g., \breaktt{src/knowledge\_graph/},
\breaktt{src/visualization/}) is governed by dedicated \texttt{SKILL.md}
files serving as Model Context Protocol (MCP) prompt-boundaries. These
explicitly define the ``rules of engagement'' for autonomous AI
inference agents---such as enforcing the zero-mock testing philosophy
via local HTTP proxies, handling specific LLM fallback parsing logic,
and respecting headless rendering constraints. This design ensures that
future AI orchestrators can natively interface with, scale, and refine
the computational meta-analysis pipeline safely and deterministically
without structural micromanagement.
\end{block}

\begin{block}{The Discovery Engine and Future Architectures}
\protect\phantomsection\label{the-discovery-engine-and-future-architectures}
Broadening our synthesis of knowledge graphs and LLMs, future iterations
of this pipeline may interface with architectures like the
\emph{Discovery Engine} \citep{baulin2025discovery}. This comprehensive
framework is designed to overcome the limitations of the
document-centric publishing paradigm by transforming unstructured
scientific literature into a machine-operable ``world model.'' Their
approach uses systematic, self-consistent LLM distillation to extract
typed ``knowledge artifacts'' from publications, which are sequentially
assembled into a hierarchical Conceptual Nexus Model (CNM) graph and
encoded as a high-dimensional Conceptual Nexus Tensor. By explicitly
modeling experimental variables, causal relations, and evidential
contradictions within a FAIR-aligned representation, this architecture
enables AI agents to mathematically navigate the knowledge landscape,
trace provenance, and generate novel hypotheses through operations akin
to tensor factorization and Vector Symbolic Architectures (VSA). This
shift from static digital libraries to a computable, relation-rich
evidence graph deeply parallels our objective of translating
unstructured Active Inference literature into a quantifiable assertion
tracking system.
\end{block}
\end{frame}

\begin{frame}[allowframebreaks]{Open Questions}
\protect\phantomsection\label{open-questions}
This meta-analysis surfaces four empirically testable questions whose
answers would directly advance the four-step research agenda outlined in
\S\ref{sec:next_steps}.

\begin{itemize}
\item
  \textbf{Recency bias in citation weighting (Methodological
  limitation).} The citation-weighting function
  \(w(a) = \log(1+\text{citations}) \cdot \text{confidence}\)
  systematically underweights recent papers (2024--2026) which have few
  citations. A 2024 paper with 1 citation is weighted approximately
  0.69\times versus a 2015 paper with 100 citations at approximately
  4.6\times. Future work may explore time-decay normalization to
  mitigate this recency penalty.
\item
  \textbf{Domain classifier over-assignment to A2 (Philosophy).} The
  keyword-based domain classifier tends to over-assign papers to the
  broad A2 (philosophy) category, where FEP universality is implicitly
  invoked but rarely explicitly tested. This classification bias likely
  inflates H1's neutral evidence count and should be addressed in future
  work through embedding-based classification or expert annotation.
\item
  \textbf{Classifier calibration (feeds Step 1).} What proportion of A1
  (Formal Theory) papers would be reclassified under an embedding-based
  or expert-annotated scheme, and how does this affect the field's
  theoretical core? An embedding-classifier trained on a 200-paper
  labeled set and evaluated on held-out A1 vs.~A2 examples would
  quantify the fraction of ``philosophy'' papers that carry formal
  mathematical content, directly sharpening both retrieval scope and
  outcome-grounding rate.
\item
  \textbf{Falsifiability and explicit testing (feeds Step 3).} H1 (FEP
  Universality) produces a predominantly neutral evidence profile,
  consistent with the critique that FEP accommodates any behavior
  without generating distinctive predictions \citep{colombo2021free}.
  Can hypothesis definitions---and author reporting standards---be
  reformulated to require a formal, refutable empirical prediction
  before contributing a supporting assertion? The proposed
  outcome-indicator taxonomy (§\ref{sec:next_steps}) would
  operationalize this: only assertions paired with a measurable outcome
  indicator would count as empirical support, converting the neutral H1
  tally into a decomposed ``invoked vs.~tested'' breakdown.
\item
  \textbf{The Scalability Gap (feeds Step 3).} H5 (AIF Scalability)
  shows a strong positive trend, yet head-to-head comparisons with deep
  RL remain concentrated on a narrow set of benchmarks (predominantly
  low-dimensional discrete environments). Beyond what state-space
  dimensionality and reward density does the expected-free-energy
  exploration advantage of model-based AIF degrade relative to
  model-free architectures such as SAC or PPO? Answering this requires
  assembling the outcome-indicator-tagged evidence (Step 3) and
  identifying which benchmark comparisons are already in the literature
  versus which are genuinely absent.
\item
  \textbf{Evidence Cross-Pollination (feeds Step 1 + Step 4).} To what
  extent do mathematical structures underlying variational free energy
  minimization and energy function optimization in Energy-Based Models
  (VAEs, contrastive divergence) converge? Extending the corpus to
  include EBM literature (Step 1) and running the assertion extractor on
  the merged set would produce a cross-domain hypothesis score for the
  shared-architecture claim---a direct test of convergence rather than a
  theoretical argument. The rubric's corpus recall metric (Step 4) would
  validate whether the expanded retrieval actually captures the EBM
  literature at recall \(\geq 0.85\).
\end{itemize}
\end{frame}

\begin{frame}[allowframebreaks]{Pipeline as a Community Instrument}
\protect\phantomsection\label{pipeline-as-a-community-instrument}
The four next steps are not a private development roadmap---they are an
invitation. The repository is structured so that each step can be
contributed incrementally: a new source connector (Step 1), a revised
extraction prompt with evidence-quote fields (Step 2), a YAML file
defining outcome indicators per hypothesis (Step 3), and an annotation
script that computes rubric scores against a provided gold set (Step 4).
None of these require modifying the scoring engine or the knowledge
graph schema. By publishing the rubric thresholds alongside the current
baseline scores, this work makes explicit what it would take for a
community contributor to demonstrably improve the system---and provides
the tooling to verify that improvement without relying on subjective
assessment.
\end{frame}

\begin{frame}[allowframebreaks]{Limitations}
\protect\phantomsection\label{limitations}
Recency bias: The citation-weighting function
\(w(a) = \log(1+\text{citations})\cdot\text{confidence}\) systematically
underweights recent papers (2024--2026) which have few citations. A 2024
paper with 1 citation is weighted \(\sim0.69\times\) versus a 2015 paper
with 100 citations at \(\sim4.6\times\). Future work may explore
time-decay normalization.

Classifier bias: The assertion counts are also sensitive to corpus
composition: H1's large neutral tally (703) partially reflects the
keyword classifier's tendency to assign papers to the broad A2
(philosophy) category, where FEP universality is implicitly invoked but
rarely explicitly tested. This classifier bias likely inflates H1's
neutral classification count and should be addressed in future work.
\end{frame}

\end{document}
